\section{Related Work}
\label{sec:related}

\paragraph{Generative recommendation and semantic item IDs.}
Recent generative recommenders reformulate recommendation as constrained sequence generation over discrete item tokens. TIGER introduced generative retrieval for recommendation and showed that autoregressive decoding over item codes can replace traditional ranking \citep{rajput2023tiger}. LCRec aligned collaborative semantics with language-model adaptation and demonstrated that tokenized item representations can be integrated into large language model pipelines \citep{zheng2023lcrec}. MiniOneRec pushed this line further by releasing an end-to-end open-source framework that covers semantic ID construction, SFT, RL, and evaluation \citep{kong2025minionerec}. These works establish the importance of tokenizer quality, but they mostly inherit a semantic-first tokenizer view. Our work stays within the MiniOneRec pipeline and asks a narrower question: how should collaborative structure enter a \emph{hierarchical} semantic tokenizer without destabilizing it?

\paragraph{Collaborative-aware tokenization for generative recommenders.}
Several recent papers argue that text-only tokenization is insufficient. ReSID motivates recommendation-native tokenizer design from the perspective of prefix-conditional predictability and semantic quantization beyond plain language modeling \citep{liang2026resid}. PRISM introduces purified collaborative representations and semantic anchoring to avoid noisy collaborative tokenization \citep{fang2026prism}. PIT further highlights the volatility of collaborative signals and studies dynamic personalized tokenization \citep{wang2026pit}. DiscRec separates semantic and collaborative factors before fusion \citep{liu2025discrec}. These methods all strengthen the case that collaboration matters, but they still leave open the question that matters most for hierarchical IDs: should the same collaborative signal be reused across every SID level, or should different levels preserve different structures? MGR-SID is built precisely around that missing allocation question.

\paragraph{Tokenization objectives beyond reconstruction.}
Another relevant thread asks tokenizers to preserve recommendation structure rather than optimizing reconstruction alone. ETEGRec argues that tokenization should be aligned with the recommender objective itself instead of being treated as an isolated preprocessor \citep{liu2024etegrec}. CoST uses contrastive quantization to preserve semantic tokenization quality for generative recommendation \citep{zhu2024cost}. These works are close in spirit to our method because they move the burden of recommendation structure into tokenizer learning. Our difference is that we use graph-structured collaborative signals as explicit structural targets and distribute them across SID levels, instead of introducing a single tokenization objective that treats all relations uniformly.

\paragraph{Graph signal processing and transition-aware recommendation.}
The graph view bank in MGR-SID is inspired by work showing that different graph scales capture different recommendation patterns. GSPRec decomposes graph signals into different spectral components for recommendation \citep{rabiah2025gsprec}, while FaGSP studies frequency-aware graph signal processing for collaborative filtering \citep{xia2024fagsp}. Collaboration and Transition emphasizes that long-range collaborative structure and short-range transition structure are complementary rather than interchangeable \citep{zhu2023collaborationtransition}. These papers are not SID tokenizers, but they provide a strong clue for our design: the collaborative information most useful for a hierarchical tokenizer is likely multi-resolution, and middle-scale graph structure deserves explicit modeling rather than being approximated by a crude interpolation between global and local views.

\paragraph{Positioning of this work.}
MGR-SID combines ideas from all three threads above, but it does not aim to compete as a generic graph encoder or a full end-to-end personalized tokenizer. The current v1 is deliberately minimal. It inherits MiniOneRec's semantic \rqvae{} backbone, borrows the denoising discipline of PRISM, uses the predictability-based motivation of ReSID, and instantiates graph-aware supervision with a coarse/mid/local graph bank whose middle view is selected by empirical probes inspired by GSPRec and FaGSP. In this sense, the paper's main point is not ``better graph features,'' but \emph{level-specific graph supervision for semantic ID learning}.
